\section{Small curvature and strict ordinary types}
\label{sec:small-curvature-strict-types}

Multiplier ideals can miss an unbounded residual. To decide whether an
ordinary type containing such a residual admits a strict current, that
is, a current $\theta+\ddc v$ bounded below by a positive multiple of a
Hermitian form, we need representatives in that same ordinary type with
a controlled curvature bound. Fix a Hermitian form $h$ on $X$ and a
quasi-plurisubharmonic function $q$ such that
\begin{equation}\label{eq:small-curvature-type}
 \text{for every }\delta>0\text{ there is }q_\delta\sim q
 \quad\text{with}\quad\ddc q_\delta\geq-\delta h.
\end{equation}
This property does not depend on the choice of $h$, since any two
Hermitian forms are uniformly comparable on $X$.

\begin{lemma}\label{lem:small-curvature-zero-lelong}\index{Singularity type!small negative curvature}\index{Polar curvature!vanishing for small curvature}
If \eqref{eq:small-curvature-type} holds, then $q$ has zero point Lelong
numbers and zero polar curvature. The pullback of $q$ to every smooth
connected compact modification has the same small-curvature property,
and its point Lelong numbers and polar curvature also vanish.
\end{lemma}
\begin{proof}
The function $q_\delta/\delta$ belongs to $\PSH(X,h)$.
The uniform Lelong bound of \cref{lem:gauduchon-uniform-lelong}, and
invariance under bounded changes, give
\[
 0\leq\nu(q,x)=\nu(q_\delta,x)
       =\delta\nu(q_\delta/\delta,x)\leq C_h\delta.
\]
Letting $\delta$ tend to zero proves the first assertion.

For polar curvature, put $S=\{q=-\infty\}$ and let
$R=\mathbf1_S\ddc q$, a closed positive current by
\cref{prop:polar-curvature-closedness}. Choose a Gauduchon metric $\gamma$
when $n\geq2$, using \cref{thm:gauduchon}; when $n=1$ take any Hermitian
form and use $\gamma^0=1$. Bounded-type invariance in
\cref{prop:bounded-type-polar-invariance} gives
$R=\mathbf1_S(\delta h+\ddc q_\delta)$. The positive current
$\delta h+\ddc q_\delta$ has total $\gamma$-trace
$\delta\int_Xh\wedge\gamma^{n-1}$ by distributional Stokes. Consequently
\[
 0\leq\int_XR\wedge\gamma^{n-1}
      \leq\delta\int_Xh\wedge\gamma^{n-1}.
\]
Thus $R=0$ by positivity and $\delta\downarrow0$.

Now let $\pi\colon X'\to X$ be a smooth connected compact modification,
and fix a Hermitian form $h'$ on $X'$. The pullbacks $\pi^*q=q\circ\pi$
and $\pi^*q_\delta$ are quasi-plurisubharmonic by
\cref{prop:psh-descent-modification}. Their ordinary inequalities hold
pointwise, including at poles, and $\pi^*h\leq Ch'$ for some fixed $C$.
Hence $\pi^*q_\delta\sim\pi^*q$ and
$\ddc\pi^*q_\delta\geq-C\delta h'$. Since $C\delta$ can be arbitrarily
small, the preceding argument on $X'$ proves both vanishing assertions.
\end{proof}

The point Lelong statement and the polar statement use different parts
of the current: an isolated logarithmic pole in dimension at least two
has a positive Lelong number, but the complex Hessian of its potential
gives it no mass. Small curvature controls both.

\begin{lemma}\label{lem:zero-lelong-strict-smoothing}\index{Regularization!strict current with zero Lelong numbers}
Let $B$ be a smooth real $(1,1)$-form. If $q'$ is quasi-plurisubharmonic,
all its point Lelong numbers vanish, and $B+\ddc q'\geq\epsilon h$ with
$\epsilon>0$, then there is a smooth real $f$ such that
$B+\ddc f\geq\epsilon h/2$.
\end{lemma}
\begin{proof}
Apply the smooth regularization of closed currents recalled in
\cref{subsec:psh-closed-regularization} to $T=\ddc q'$, with smooth
representative zero and the lower bound $\gamma=\epsilon h-B$, which
need not be closed. Since $X$ is compact, the Hermitian form used there
is bounded by a constant multiple of $h$. The regularization therefore
gives smooth functions $f_k$ on $X$ and a constant $C$ with
\[
 \ddc f_k\geq\gamma-C\lambda_kh-\epsilon_kh,
\]
where the continuous nonnegative functions $\lambda_k$ decrease
pointwise to the Lelong numbers of $q'$, which vanish, and
$\epsilon_k\downarrow0$. Dini's theorem gives uniform convergence of
$\lambda_k$ to zero. Choose $k$ with
$C\|\lambda_k\|_\infty+\epsilon_k\leq\epsilon/2$ and set $f=f_k$;
then $B+\ddc f\geq\epsilon h/2$.
The smooth approximants need not remain in the ordinary type of $q'$.
\end{proof}

Let $\Delta=\sum_{j=1}^Na_jD_j$ be a finite effective real divisor,
where the $D_j$ are distinct prime divisors and $a_j>0$. Choose smooth
Hermitian metrics on the line bundles $\mathcal O_X(D_j)$, with canonical
sections $s_j$, and put
\[
 \sigma_\Delta=\sum_ja_j\log\|s_j\|^2,\qquad
 \ddc\sigma_\Delta=[\Delta]-\chi_\Delta,
\]
where $\chi_\Delta$ is smooth and closed. The case $\Delta=0$ means
$\sigma_\Delta=\chi_\Delta=0$; zero coefficients are omitted from the
logarithmic sum.

\begin{theorem}\label{thm:small-curvature-strict-type}\index{Strict current!small-curvature residual}\index{Divisor!strict ordinary type}
Suppose \eqref{eq:small-curvature-type} holds and put
$B=\theta-\chi_\Delta$. The ordinary type of $\sigma_\Delta+q$
contains a strict $\theta$-current, that is, some $v\sim\sigma_\Delta+q$
satisfies $\theta+\ddc v\geq\epsilon h$ with $\epsilon>0$, if and only
if $B$ has a smooth strictly positive representative: $B+\ddc f>0$ for
some smooth real $f$.
\end{theorem}
\begin{proof}
Suppose first that $v\sim\sigma_\Delta+q$ and
$\theta+\ddc v\geq\epsilon h$ with $\epsilon>0$.
Lelong numbers are additive and invariant under bounded changes, so
\cref{lem:small-curvature-zero-lelong} gives $\nu(v,D_j)=a_j$.
Apply \cref{lem:nonclosed-divisorial-removal} with background
$\theta-\epsilon h$ to obtain
\[
 \theta+\ddc v-[\Delta]\geq\epsilon h.
\]
For $\Delta=0$ this is the original strict inequality.
The $L^1$ difference $v-\sigma_\Delta$ consequently has a global
quasi-plurisubharmonic representative $q'$ with
$B+\ddc q'\geq\epsilon h$. Off the divisor it satisfies $q'\sim q$;
these inequalities extend everywhere by
\cref{lem:qpsh-almost-everywhere-order}. Thus $q'$ has zero point
Lelong numbers, and \cref{lem:zero-lelong-strict-smoothing} gives
$B+\ddc f>0$ for a smooth $f$.

Conversely, compactness gives $B+\ddc f\geq\epsilon h$ for some
$\epsilon>0$. Choose $q_\delta$ in \eqref{eq:small-curvature-type} with
$\delta<\epsilon/2$. Then $v=\sigma_\Delta+q_\delta+f$ is in the
specified ordinary type, and Poincar\'e--Lelong gives
\[
 \theta+\ddc v=[\Delta]+B+\ddc f+\ddc q_\delta
                    \geq\epsilon h/2.
\]
\end{proof}
